\section{Scope and the result of the prescribed test}
This is Turn 2 of the ten-turn ES programme. Its task is to join the complete
joint exterior/middle obstruction to the universally available shifted-factor
system. The canonical system, its reciprocity laws, and existence of a cycle
are inherited from Bryan's note \cite[\S\S1--7]{BryanGraph}. We retain \emph{all}
nonresidue factors, extend Turn 1 to its auxiliary composite vertices, calculate
four transition types and their exact common domains, and give a complete
closed-component decision theorem. The theorem that every closed component
contains an occupied local selector is false. We certify a closed seven-cycle
at $p=2521$, and show a primitive sextic obstruction feeding a closed component
at $p=3361$. These are counterexamples to that proposed proof step, not to ES.
No historical-priority claim is made. The complete graph construction below is
proved directly; none of the source's status labels replaces a proof.

\section{The universal original graph, with both vertex types}
Let $p$ be prime in one of the six classes
$1,121,169,289,361,529\pmod{840}$. Thus $p\equiv1\pmod{24}$ and
$(2/p)=(3/p)=(5/p)=(7/p)=1$, by quadratic reciprocity and its supplementary laws.
Here and below parentheses denote the Legendre symbol at a prime denominator
and the Jacobi symbol at a positive odd composite denominator.
Set
\[
 V_p=\{q<p:q\text{ prime},\ (q/p)=-1\},\qquad
 \sigma_q=\begin{cases}1&q\equiv3\pmod4,\\3&q\equiv1\pmod4,\end{cases}
 \quad R_q=\sigma_q q,\quad A_q=(p+R_q)/4.                 \tag{1}
\]
\begin{lemma}[Complete source and outgoing factors]
$V_p$ is nonempty, its members are at least 11, and for every $q\in V_p$,
\[
 p/4<A_q<p,\quad R_q\equiv3\pmod4,\quad \gcd(pA_q,R_q)=1,
 \quad (A_q/p)=-1.                                      \tag{2}
\]
Write $A_q=\prod_t t^{e_{qt}}$. The complete outgoing edge set is
\[
 q\longrightarrow r\quad\Longleftrightarrow\quad
 r\mid A_q,\ r\in V_p,
 \qquad\text{with retained weight }m_{qr}=e_{qr}.          \tag{3}
\]
It is nonempty, has no self-edge, and
$\sum_{r\in V_p}m_{qr}$ is odd. Even positive multiplicities are retained.
\end{lemma}
\begin{proof}
A positive nonresidue below $p$ has a nonresidue prime factor below $p$.
The four displayed residue symbols exclude $2,3,5,7$.
Integrality and $p/4<A_q<p$ follow from (1) and $q<p$. Since a common divisor
of $A_q,R_q$ divides $p$, and $p$ divides neither $q$ nor $3$, (2)'s coprimality
holds. Modulo $p$, $4A_q=\sigma_q q$, so $(A_q/p)=-1$.
Factoring this equality of characters gives the odd multiplicity sum and a
nonempty outgoing set. Every outgoing prime is at most $A_q<p$, and $q$ itself
cannot divide $A_q$. These assertions use all original prime multiplicities.
\end{proof}

At the vertex retain the entire centered source
\[
 \mathcal B_q=\prod_{t\mid A_q}[-e_{qt},e_{qt}]_{\mathbb Z},\quad
 v_q(\beta)=\prod_t t^{\beta_t}\pmod{R_q},\quad
 u_q(\beta)=\prod_t t^{e_{qt}+\beta_t}.                    \tag{4}
\]
The last map is a bijection to $\operatorname{Div}(A_q^2)$, inverse
$\beta_t=v_t(u)-e_{qt}$. Keep the integer function
$\mu_q(g)=\#v_q^{-1}(g)$, its support $W_q$, and separately the canonical hits
\[
 E_q=\{u\mid A_q^2:R_q\mid4u+1\},\qquad
 M_q=\{u\mid A_q^2:R_q\mid u+A_q\}.                      \tag{5}
\]
The targets on (4) are $-p^{-1}$ and $-1$. Inversion of the \emph{same} exponent
vector gives the reflected exterior target $-p$, not a third channel.

\paragraph{Marked return at an auxiliary vertex.}
For a hit put $a=A_q,R=R_q,d=\gcd(a,u)$ and
\[
 (h,r_0,s_0)=(d^2/u,u/d,a/d),\quad a=hr_0s_0,\quad u=hr_0^2.
\]
Prime valuations prove integrality and $\gcd(r_0,s_0)=1$. The ordered returns are
\[
 \begin{array}{ll}
 E:&\kappa=(pr_0+s_0)/R,\quad (x,y,z)=(a,hs_0\kappa,phr_0\kappa),\\
 M:&\lambda=(r_0+s_0)/R,\quad (x,y,z)=(a,phs_0\lambda,phr_0\lambda).
 \end{array}                                               \tag{6}
\]
The gates and units imply the quotients are integers; multiplying the reciprocal
sum by $phr_0s_0\kappa$ or $phr_0s_0\lambda$ proves $4/p$.
The ordered inverse is
\[
 u=\frac{a^2}{Ry-pa}\ (E),\qquad
 u=\frac{pa^2}{Ry-pa}\ (M).                               \tag{7}
\]
These are inherited Type I/II coordinates \cite[\S2, Propositions 2.2, 2.6]{ET}.
The preserved record is
\[
(p,q,\sigma_q,R_q,A_q,u,\beta,E/M,h,r_0,s_0,\kappa\text{ or }\lambda;x,y,z).
\]

The prime-residual vertices satisfy $A_q<p/2$. The composite vertices need not.
If $A_q>p/2$, there can be no M hit: its last two denominators are multiples of
$p$, and the reciprocal sum would be at most $1/A_q+2/p<4/p$.
For an E hit in this larger range, $z\ge p$ and $y<p/2$. To verify the latter,
order any ES triple $x\le y\le z$. If $x>p/2$, then $y<p$, $z=pk$, and
$t=(4k-1)/p$ is a positive integer $3\pmod4$. Writing $x=gA,y=gB$,
$\gcd(A,B)=1$, gives $k=ABD$, $tg=D(A+B)$, $\gcd(t,D)=1$.
The inequality $x>p/2$ implies $2DA(A-B)>-1$, so $A\ge B$; ordering gives
$A=B=1$, whence $t\mid2$, a contradiction. Equality is impossible for odd $p$.
The coordinate switch is explicitly
$(h,r_0,s_0,\kappa)\mapsto(h,\kappa,s_0,r_0)$.
Indeed $\gcd(\kappa,s_0)=1$ follows from $R\kappa=pr_0+s_0$ and the units;
also $(4hs_0\kappa-p)r_0=p\kappa+s_0$.
Thus its divisor is $u'=h\kappa^2$, and the switch is its own inverse.
Swap the two non-$p$ denominators, set $a'=y,R'=4y-p$ and
$u'=y^2/(R'a-py)$, and apply (6)--(7). The original permutation $(y,a,z)$ is
retained. This is a proven return to the first-half atlas, not deletion of a vertex.

\section{The Turn 1 obstruction attaches without changing the source}
Define, independently at every vertex,
\[
 G_q=(\mathbb Z/R_q\mathbb Z)^\times,\quad
 \Gamma_q=\langle-1,2,t:t\mid A_q\rangle,\quad
 K_q=\operatorname{Stab}_{G_q}(W_q),\quad\mathcal Q_q=\Gamma_q/K_q.
                                                               \tag{8}
\]
Since $1\in W_q$, $K_q\subseteq W_q\subseteq\Gamma_q$.
With $\alpha=-1K_q,\zeta=2K_q,\eta=pK_q,g_t=tK_q$,
\[
 \eta=\zeta^2\prod_tg_t^{e_{qt}},\qquad
 F_q=\{\alpha,\alpha\eta,\alpha\eta^{-1}\}.                \tag{9}
\]
The distinct target count is $1,2,3$ according as $\eta=1$, $\eta\ne1$ is an
involution, or $\eta^2\ne1$. All three labels are retained. Joint failure is
exactly $(W_q/K_q)\cap F_q=\varnothing$.

All the algebraic results of Turn 1 apply on (2), including composite residuals
and $A_q>p/2$: their proofs use only the unit condition and $4A_q=p+R_q$,
not the first-half size bound. For completeness their principal restrictions are
\[
 2\sum_{g_t\ne1}e_{qt}\le|\mathcal Q_q|-1-|F_q|,
 \qquad \operatorname{ord}(g_t)>2e_{qt}+1\ (g_t\ne1).       \tag{10}
\]
For any nonempty subset $I$ of active prime indices, use the actual remaining
centered word set $U_I$ and $C_I=\langle g_t:t\in I\rangle$; then
\[
 2\sum_{t\in I}e_{qt}
 \le |C_I|-1-\max\{1,|(F_qU_I^{-1})\cap C_I|\}.          \tag{11}
\]
Indeed any period of a factor subproduct persists to the final quotient support.
Each of its $2e$ nonidentity two-point factors must strictly enlarge the support;
a full cyclic interval would create a period. Comparing this growth with the
actual pulled-back missing targets proves (10)--(11).
The complete order-at-most-eight classification and its fine coefficients are
retained from Turn 1 \cite[\S\S3--6]{Turn1}, rather than reclassified here.
In particular, effective index six at a prime vertex forces exactly one simple
nonresidue factor $r$, with
\[
 A_q=Br,\quad W(B)=K_q,\quad |rK_q|=6,\quad
 pK_q=rK_q\text{ or }r^{-1}K_q.                          \tag{12}
\]
The general prime-shift sextic antecedent is \cite{BryanSix}; the precise
saturated inactive source is in Turn 1.

Each quotient has its own section $s_q$ and carry
$\omega_q(c,d)=s_q(c)s_q(d)s_q(cd)^{-1}\in K_q$.
The retained fine coefficient is $\mu_q(s_q(c)k)$, not a uniform share of its
coset mass. The new graph does not identify $K_q$ and $K_r$, their sections,
or their quotient coefficient arrays merely because an edge exists.

\section{All four arithmetic transition types}
An edge $q\to r$ retains
\[
 m=v_r(A_q),\quad B=A_q/r^m,\quad c=A_q/r=r^{m-1}B,
 \quad\boxed{p+\sigma_q q=4cr.}                          \tag{13}
\]
These are bijective coordinates for the edge after its input prime and complete
factorization are retained. If $\tau=\sigma_r$, the next source is exactly
\[
 \boxed{4A_r=(4c+\tau)r-\sigma_q q,\qquad R_r=\tau r.}   \tag{14}
\]
No factorization of $A_r$ is inferred from that of $A_q$; it is a new original
factorization of this specified integer.

For every prime $t\mid A_q$ the common character identity is
\[
 \boxed{(t/R_q)=(t/p).}                                  \tag{15}
\]
For odd $t$, use $R_q\equiv-p\pmod t$, $R_q\equiv3\pmod4$,
and quadratic reciprocity. The sign from $(-1/t)$ cancels the reciprocity sign.
For $t=2$, its occurrence makes $R_q\equiv-p\pmod8$, giving the same formula.
For an outgoing edge this yields the already inherited source-prime rules
\[
 (r/q)=\begin{cases}-1&\sigma_q=1,\\-(-3/r)&\sigma_q=3,\end{cases}
 \qquad
 (q/r)=\begin{cases}-(-1/r)&\sigma_q=1,\\-(-3/r)&\sigma_q=3.
 \end{cases}                                               \tag{16}
\]
The reversed symbol need not be a nonresidue. In particular, an all-$3\pmod4$
two-cycle is impossible. A mixed two-cycle requires its $3\pmod4$ prime to be
$1\pmod3$, and a two-cycle with both primes $1\pmod4$ requires equal residues
modulo three. These follow by applying both directions of (16) and reciprocity;
they are necessary restrictions, not a converse existence theorem.

The fine inverse-exterior target at the successor is always $-p\pmod{R_r}$.
Its complete transfer table is
\[
\begin{array}{c|c|l}
 \sigma_q&\tau&\text{retained equality at the successor}\\\hline
1&1&q\equiv-p\pmod r\\
3&1&3q\equiv-p\pmod r\\
1&3&q-4cr\equiv-p\pmod{3r}\\
3&3&3q-4cr\equiv-p\pmod{3r}.
\end{array}                                                \tag{17}
\]
For $\tau=3$ the exact CRT representation is
\[
 \boxed{-p\longleftrightarrow(2,\sigma_q q\bmod r)
 \in(\mathbb Z/3\mathbb Z)^\times\times(\mathbb Z/r\mathbb Z)^\times.}
                                                               \tag{18}
\]
The inverse sends $(a,b)$ to
\[
 a r(r^{-1}\bmod3)+3b(3^{-1}\bmod r)\pmod{3r}.
\]
In (18), $a=2$. In the last row of (17), $3q$ itself is a nonunit modulo
$3r$. The cofactor correction must be retained; division by three in that
unit group is undefined. The three local targets all have first CRT coordinate
2, and their second coordinates are $-1,-p,-p^{-1}\pmod r$. The full joint
coefficient is obtained from the \emph{same} exponent vector in both coordinates.
Independently chosen local witnesses are not a global witness.

\paragraph{The exact common original-divisor domain.}
Put $\Delta=A_r-A_q=(\tau r-\sigma_q q)/4$. Then
\[
 \gcd(R_q,R_r)=\begin{cases}3&\sigma_q=\tau=3,\\1&\text{otherwise},\end{cases}
 \quad
 \operatorname{Div}(A_q^2)\cap\operatorname{Div}(A_r^2)
 =\operatorname{Div}(\gcd(A_q,A_r)^2).                    \tag{19}
\]
The second assertion is primewise minimum of the two original exponent bounds.
On this precise common domain, fixed $u$ transports its centered exponents by
\[
 \beta_t^{(r)}=\beta_t^{(q)}+v_t(A_q)-v_t(A_r),
 \quad g_E'=g_E=4u+1,\quad g_M'=g_M+\Delta.              \tag{20}
\]
The maps (19)--(20) have their literal inclusion inverses, not a map from all
source words to all target words. Even on their common domain the gate can change.
At $p=1801$, $19\to13$ has $A_{19}=455$, $A_{13}=460$; their common divisors
of squares are $1,5,25$. The original $u=1$ M hit at modulus 19 is a miss at
modulus 39. Its source denominators are $(455,19666920,43224)$.

If two successive prime vertices both have primitive index six, (17) puts the
previous vertex in an order-three forbidden class at the successor, while the
new exceptional prime has order six. With actual effective quotient $K_r$,
$qK_r=t^{2\varepsilon}K_r$, $\varepsilon=\pm1$, for the next active prime $t$.
Only when the ambient index is six may this be written as a sixth-power congruence.
The earlier consecutive example is attributed to \cite[\S\S2--4]{BryanChain}.

\section{A complete closed-failure test}
Let $\mathscr H=\{q:E_q\cup M_q\ne\varnothing\}$ and let $N^+(q)$ be the
complete distinct successor set, with multiplicities kept separately.
Define
\[
 U_0=V_p\setminus\mathscr H,\qquad
 U_{j+1}=\{q\in U_j:N^+(q)\subseteq U_j\}.               \tag{21}
\]
\begin{theorem}[Finite closure and exact reachability certificates]\label{thm:closure}
The sequence stabilizes in at most $|V_p|$ iterations. Its terminal set $U_\infty$
is the unique largest forward-closed set of locally failed vertices. For every
original vertex,
\[
 \boxed{q\notin U_\infty\iff
 \text{some path of length at most }|V_p|-1\text{ reaches }\mathscr H.}
                                                               \tag{22}
\]
Moreover $U_\infty\ne\varnothing$ if and only if the complete directed graph
has a sink strongly connected component disjoint from $\mathscr H$.
\end{theorem}
\begin{proof}
Each nonstationary step removes a vertex, and the final set is closed and
failed by its definition. Every closed failed set is contained in every $U_j$,
by induction. A removed vertex either is a hit or has a successor removed
earlier, producing a strictly decreasing removal rank and a terminating hit
path. Conversely a shortest hit path has no repeated vertex, so has length at
most $|V_p|-1$; backwards induction removes each of its vertices.
A nonempty finite closed set has a bottom strongly connected component. Because
the set is closed, that component has no outgoing edge in the full graph.
Conversely any failed sink component remains in every $U_j$.
\end{proof}

The executable inverse of (22) records an outgoing edge decreasing the shortest
hit distance and terminates in an actual state (6). A certificate of nonreachability
contains its entire successor-closed failed set. A sink component can be certified
by all its outgoing factorizations and two rooted reachability trees within it.
A selected failed cycle with an omitted exit is not such a certificate.
The graph-theoretic SCC decomposition is standard \cite[Theorems 3.10 and 5.2]{Tarjan}; (21)--(22)
are proved here with their arithmetic observation and full-edge semantics.

The test is complete at fixed $p$: a sieve supplies all primes below $p$;
trial division or a smallest-factor table factors each $A_q<p$; the complete
boxes (4) decide both channels. There are fewer than $p$ vertices, and each
has at most $2A_q-1<2p$ square-divisor candidates, by divisor pairing. Thus fewer
than $2p^2$ candidates and twice that many gate tests suffice as a crude finite
bound. No favourable factorization, density claim, or supplied no-hit flag
is assumed. The graph remains a specified selector subfamily, not the complete
ES atlas: $U_\infty=V_p$ would not by itself certify an ES counterexample.

\section{Complete cyclic integer coordinates}
For a simple directed cycle $q_0,\ldots,q_{l-1}$, retain
$p+\sigma_iq_i=4c_iq_{i+1}$ with indices modulo $l$. Set
\[
 D=4^l\prod_i c_i,\qquad S=\prod_i\sigma_i.
\]
At each cyclic starting point define $T_i$ by iterating
\[
 D_0=1,\ S_0=1,\ T_0=0,\qquad
 D_{j+1}=4c_{i+j}D_j,\quad S_{j+1}=\sigma_{i+j}S_j,
 \quad T_{j+1}=\sigma_{i+j}T_j+D_j.                       \tag{23}
\]
Let $T_i$ denote the final $T_l$ for that starting point.
\begin{theorem}[Integral cycle inverse]
Every such cycle satisfies
\[
 \boxed{(D-S)q_i=pT_i.}                                   \tag{24}
\]
In particular, with $h=\gcd(T_0,\ldots,T_{l-1})$,
\[
 \boxed{p=(D-S)/h,\qquad q_i=T_i/h.}                       \tag{25}
\]
Thus an ordered coefficient word $(\sigma_i,c_i)$ encodes at most one canonical
prime cycle. Its converse is a finite exact test: calculate (23)--(25), then
verify integrality, distinct prime vertices below the prime $p$, the specified
$\sigma$ types, nonresidue symbols, factor edges and, when claimed, all-edge
closure and both local failures. None of those latter tests is omitted.
\end{theorem}
\begin{proof}
Composition of the affine equations yields
$D_jq_{i+j}=S_jq_i+pT_j$, proving (23) and then (24).
$D>S$ because every $4c_i>\sigma_i$. Since $p$ and $q_i$ are distinct primes,
$p\mid D-S$. Put $h_0=(D-S)/p$; then $T_i=h_0q_i$.
A simple cycle has at least two distinct primes, so their common gcd is one;
hence $h=h_0$. This proves (25) and its uniqueness statement.
\end{proof}
Forgetting the cycle's starting point loses its $l$ rotations. They are not
silently identified in (23). The affine maps have slopes at most $3/4$ and
consistent positive fixed points; their contraction is not a contradiction to
cycle existence. The next example satisfies (24) exactly.

\section{A closed all-failure component actually exists}
\begin{theorem}[Certified obstruction to canonical-cycle closure]\label{thm:counter}
At the hard prime $p=2521$, the complete canonical graph has the sink component
\[
 \boxed{11\to211\to683\to89\to17\to643\to113\to11.}     \tag{26}
\]
Every vertex has exactly its displayed outgoing nonresidue prime, and every
one of its two local selectors is empty. It is therefore impossible to reach
an occupied canonical vertex from any point of (26). This refutes the assertion
that every closed component has an occupied selector, and refutes such a
reachability assertion starting at the least nonresidue prime.
\end{theorem}
\begin{proof}
All displayed integers are prime, checked by complete trial division, as is
$p$. The exact source data are
\[
\begin{array}{r|r|l|r|r|r}
 q&R_q&A_q\text{ with factorization}&r&m_{qr}&|\Gamma_q/K_q|\\\hline
11&11&633=3\cdot211&211&1&10\\
211&211&683&683&1&210\\
683&683&801=3^2\cdot89&89&1&682\\
89&267&697=17\cdot41&17&1&176\\
17&51&643&643&1&32\\
643&643&791=7\cdot113&113&1&642\\
113&339&715=5\cdot11\cdot13&11&1&224.
\end{array}                                                \tag{27}
\]
Every $K_q$ here is trivial. The nonresidue membership and absence of other
outgoing factors follow either by exact Euler powers or reciprocity on the
complete factorizations. The complete centered supports and three targets are
listed in Appendix~\ref{app:cycle}, where direct comparison proves every miss.
Thus (26) is strongly connected and has no unrecorded exit.
\end{proof}
The total original divisor mass on the seven vertices is 75, not 71 (the latter
is the sum of their support sizes). The full $p=2521$ graph has 187 vertices,
four occupied vertices, and $|U_\infty|=107$. Its unique bottom component is
(26). Its occupied vertices are $23,29,37,1583$; none is reachable from this component.
For its ordered coefficient word,
\[
 (c_i)=(3,1,9,41,1,7,65),\quad(\sigma_i)=(1,1,1,3,3,1,3),
\]
\[
 D=8252375040,\quad S=27,\quad h=3273453,
 \quad D-S=2521h.                                        \tag{28}
\]
All rotated numerators $T_i=hq_i$ are retained in JSON.

The prime $2521$ itself is solved, for example, by the original middle state
$(R,a,u,h,r_0,s_0,\lambda)=(23,636,8,2,2,159,7)$:
\[
 \frac4{2521}=\frac1{636}+\frac1{5611746}+\frac1{70588}.    \tag{29}
\]
Thus the local-cycle implication is false; ES is not contradicted.

At $p=3361$ there is another closed failed component
\[
 31\to53\to11\to281\to1051\to1103\to31,                 \tag{30}
\]
with source factors $848=2^4\cdot53$, $880=2^4\cdot5\cdot11$,
$843=3\cdot281$, $1051$, $1103$, and $1116=2^2 3^2\cdot31$.
The effective indices are respectively $6,104,10,560,1050,1102$.
A genuine primitive sextic defect therefore feeds a closed component: an
increase of quotient index does not imply escape or a well-founded decrease.
Its cycle identity is $566231040-9=3361\cdot168471$.
The full graph has 233 vertices, three occupied vertices, and 191 trapped ones.

The inherited consecutive sextic chain at $p=808369$ \cite[\S4]{BryanChain}
continues to the closed failed six-cycle
\[
 19\to28871\to6977\to8293\to2003\to67531\to19.           \tag{31}
\]
Every listed node, factor, complete gate set and stabilizer is in the certificate.
This verifies the listed component only, not a fresh complete graph census at
$808369$.

\section{A legitimate same-prime extension, with its limitation}
A different multiplier can be included only as a separately marked source.
For $q\in V_p$, let $s>0$ obey
\[
 sq\equiv3\pmod4,\quad sq<3p,\quad(s/p)=1.
\]
Then $A_{q,s}=(p+sq)/4$ satisfies all source conditions (2), and its complete
nonresidue successors and two gates are defined in exactly the same manner.
For two such multipliers $s,t$,
\[
 A_{q,t}-A_{q,s}=q(t-s)/4,\quad
 \boxed{\gcd(A_{q,s},A_{q,t})=\gcd(A_{q,s},(t-s)/4).}      \tag{32}
\]
This follows from $q\nmid A_{q,s}$. For consecutive multipliers differing by
four the original factors are coprime. Their divisor sources cannot be identified.

At $p=2521,q=11$, replacing the canonical multiplier 1 by the actual quadratic
residue 5 yields $R=55,A=644=2^2\cdot7\cdot23$. Its $u=16$ M state is
\[
 (h,r_0,s_0,\lambda)=(1,4,161,3),\qquad
 \boxed{\frac4{2521}=\frac1{644}+\frac1{1217643}+\frac1{30252}.} \tag{33}
\]
The old and new $A$ values are coprime; this is a newly available source, not
transport of a missed divisor through a quotient. It repairs this component
once that extra marked port is allowed. No claim that multiplier 5, the pair
of ports $\sigma,\sigma+4$, or all permitted ports succeed universally is made.

Two further regression distinctions are useful. At $p=3361,q=977$,
$A_q=11^2\cdot13$ has both outgoing edges $11$ (weight two) and $13$ (weight one).
The former enters (30); the latter is immediately occupied. Choosing the least
factor discards a real escape. At $p=1201,q=461$ the auxiliary E hit is
\[
 (R,A,u)=(1383,646,24548),\quad(x,y,z)=(646,561,25602918).
\]
Its permutation returns the first-half E state
$(R',A',u')=(1043,561,18513)$ with the ordered triple $(561,646,25602918)$.
Deleting vertices with $A>p/2$ would delete this genuine original return.

\section{Verified scope and Turn 3 handoff}
The full scan is over all 137 hard-class primes at most 50,000. It includes
188,255 original nonresidue-prime vertices, 191,789 distinct factor edges and
193,997 edge occurrences. The complete source enumeration has 8,357,211
divisor candidates, 1,729 E states and 2,728 M states in 1,709 occupied vertices.
There are 57,416 auxiliary vertices above $p/2$, carrying ten E states; their
permutations are verified. There are 39 primes with a nonempty failure trap,
and 11,954 trapped vertices altogether. Eight least-prime starting vertices
are trapped. Every prime in this finite scan has some occupied graph vertex;
that is a finite result, not a universal theorem.

The first hard prime with a canonical trap is 2521: both implementations
independently enumerate every earlier hard prime and its complete graph.
The first implementation verifies all factor-cut bounds of Turn 1 at every
vertex for $p\le5000$, and full fine/coset records in the declared detailed
examples; it does not claim every factor-cut recomputed at 50,000.
Separate finite graph fixtures exhaust all loop-free graphs with nonempty
successor sets on two through four vertices and every hit mask: 38,636 cases.
The main and separate implementations use different character tests, gate
reconstruction, SCC algorithms and reachability calculations. They share no
code. Normal and optimized receipts and exact scopes accompany the package.

The exit of Turn 2 is the complete original graph, its all-type arithmetic
transitions, and a terminating closed-failure test with exact counterexamples
to the proposed local closure principle. Turn 3 must not try to prove that
all canonical closed components are rescued. A sufficient global claim would
be $U_\infty\ne V_p$ for every required prime, which is much weaker than the
false claim $U_\infty=\varnothing$. Alternatively an explicitly enlarged
same-$p$ source must prove its own escape mechanism. Neither follows from
reciprocity signs, a supported-zero relation, or a finite cycle alone.

\appendix
\section{Complete seven-cycle supports}\label{app:cycle}
The following sets are the supports of the full products
$\prod_{t^e\parallel A_q}\sum_{j=-e}^e[t]^j$ modulo $R_q$.
Their original multiplicities, not only these supports, are in the JSON.
In the displayed cycle order, the target triples $(-1,-p,-p^{-1})$ are
\[
 (10,9,5),\ (210,11,96),\ (682,211,560),\ (266,149,224),
 \ (50,29,44),\ (642,51,290),\ (338,191,71).
\]
\begin{align*}
W_{11}&=\{1,2,3,4,6,7,8\},\\
W_{211}&=\{1,38,50\},\\
W_{683}&=\{1,3,9,76,89,110,118,228,238,267,307,330,485,492,617\},\\
W_{89}&=\{1,17,41,46,110,163,172,238,254\},\\
W_{17}&=\{1,28,31\},\\
W_{643}&=\{1,7,92,108,113,126,148,239,387\},\\
W_{113}&=\{1,5,11,13,19,32,37,53,55,65,68,70,142,143,160,185,206,\\
&\hspace{23mm}209,214,232,247,265,266,275,313\}.
\end{align*}
Every target is absent from its corresponding support. Every original factor
and every exponent range that produces these sets is displayed in (27).
